\chapter{Wprowadzenie}
\label{chap:pierwszy}

%(zobacz rys. \ref{fig:logo})
%\begin{figure}[ht]
%	\centering
%	\includegraphics[width=50mm]{images/pipp.png}
%	\caption{Prototyp loga platformy PIPP}
%	\caption*{Źródło: opracowanie własne.}
%    \label{fig:logo}
%\end{figure}

Poniższy tekst prezentuje projekt przedsięwzięcia e-biznesowego w postaci platformy integrującej prywatnych przewoźników. Mali lokalni przewoźnicy nie są bowiem obecnie zintegrowani w takim stopniu jak transport zbiorowy w miastach, koleje, czy komunikacja dalekobieżna, co sprawa, że planowanie podróży z ich udziałem jest znacząco mniej wygodne dla podróżnych. Projekt przewiduje stworzenie platformy pozwalającej pasażerom na wyszukiwanie i planowanie tras z możliwością łatwego zakupu cyfrowych biletów wewnątrz platformy podobnej do istniejących platform stosowanych przykładowo przez przewoźników kolejowych. Sporą bazą użytkowników cieszą się portale takie jak KOLEO, BLIKOM, Portal Pasażera, E-Podróżnik, czy wiele innych. \citep{Drewnowski_Małachowski_2018} Dodatkowo podróżni zyskają dostęp do rozkładów jazdy i śledzenia nadawanej na żywo lokalizacji pojazdów różnych przewoźników w jednym miejscu. Z drugiej strony dla firm transportowych platforma będzie źródłem wielu cennych danych analitycznych pozwalających im optymalizować trasy i kursy, a w konsekwencji zwiększać efektywność i zyski. Obecnie bowiem ze względu na rozdrobnienie rynku przewozowego takie dane trudno pozyskać.
\begin{quote}
    \textit{„Wyraźnie brakuje wyczerpujących i bogatych w aktualne dane opracowań dotyczących szerzej pojętej tematyki autobusowych przewozów pozamiejskich. Jedną z głównych tego przyczyn jest postępujące rozdrobnienie rynku przewozowego, wywołujące trudności m.in. w prostej inwentaryzacji danych służących dalszym analizom.”} (\citealt{Wolański_Mrozowski_Pieróg_2016}, s. 64)
\end{quote}
Dostęp do takich bez wątpienia cennych danych powinien się przełożyć na chęć udziału przewoźników w tym przedsięwzięciu, a dodatkowo im więcej przewoźników nawiąże współpracę z platformą, tym system będzie miał więcej do zaoferowania dzięki efektowi sieci. Tym sposobem im więcej firm nawiąże współpracę z platformą tym więcej firm będzie chciało nawiązać współpracę z platformą.

Zdecydowanym atutem tego przedsięwzięcia jest także to, że wpisuje się ono w obecne światowe trendy polityczne dotyczące redukcji śladu węglowego. Jak wiadomo globalne ocieplenie i zmiany klimatu stanowią poważne wyzwanie dla ludzkości w nadchodzących latach i temat ten jest bardzo poważnie traktowany przez rządy całego świata na najwyższych szczeblach. Jako że komunikacja zbiorowa generuje mniejszy ślad węglowy na osobę w porównaniu do transportu indywidualnego (\citealt{Szymalski_Bukowicka_2022}, tab. 4 s. 80), to ułatwienie planowania podróży zbiorowym transportem regionalnym powinno przełożyć się na większe zainteresowanie tą formą transportu, a w konsekwencji na ograniczenie szkodliwych emisji. Dodatkowym elementem, który może przekonywać polityków do wspierania projektu jest walka z wykluczeniem komunikacyjnym. Z racji tego, rządy wielu krajów, a w szczególności krajów Unii Europejskiej będą przychylnie patrzeć na wszelkie inicjatywy promujące transport zbiorowy. Z dużą pewnością sprawia to, że otoczenie prawne tego przedsięwzięcia powinno być przyjazne i stabilne, co dobrze rokuje na przyszłość.
\begin{quote}\textit{„Wraz z rosnącym znaczeniem ochrony środowiska, oszczędzenia energi i redukcji śladu węglowego, zapewnienie zrównoważonego transportu stało się kluczowym wyzwaniem.”}
    %"With the increasing prominence of environmental protection, energy conservation, and carbon reduction issues [[ 1 ]], ensuring sustainable transportation [[ 3 ]] has become a critical and challenging goal."
\citep{Chung2023-ih} \end{quote}

Kolejną zaletą tej działalności jest to, że transport zbiorowy jest bardziej atrakcyjny dla ludzi o niższych dochodach, ponieważ pasażerowie go wybierający często szukają oszczędności w porównaniu do transportu indywidualnego. \citep{Nimorakiotaki2020-bl} Z tego względu usprawnienie transportu zbiorowego na terenach mało zurbanizowanych, a więc właśnie terenach o statystycznie niższych dochodach na osboę może się okazać owocne biznesowo, ponieważ mogło by zagospodarować niezaspokojone potrzeby mieszkańców właśnie tych obszarów.

Dodatkowo przedsięwzięcie ze względu na to, że skupia się głównie na oprogramowaniu jest łatwo skalowalne, co na pewno będzie atrakcyjne dla potencjalnych inwestorów. Głównym kosztem będzie początkowy wydatek na stworzenie całego systemu. Później już stosunkowo niewielkie koszty będą potrzebne do jego utrzymania. Dystrybucja cyfrowa samego oprogramowania generuje już natomiast w praktyce pomijalne koszty.